\subsection{Validation Rules} \label{subsec:validation-rules}
\begin{figure*}
\begin{center}
    % $\inference
    %     [Assignment]
    %     {\arbitrary{v} & \inpwe{\lab}{v} x \gennedby e, \out}
    %     {\inp xs \gennedby (\lab \Longleftarrow \mathit{True}) \mapsto e, \out}$
    $\inference
        [\ValidateAssignment]
        {\arbitrary{v} & \inpwe{\lab}{v} x \gennedby e, \out}
        {\inp xs \gennedby (\lab \Longleftarrow \mathit{True}) \mapsto e, \out}$

    % list operations
    $\inference
        [\ValidateCons]
        {\inp x \gennedby r_1, \out & \dinp xs \gennedby r_2, \dout}
        {\inp x \ccons xs \gennedby r_1 \cons r_2, \dout}$
    $\inference
        [\ValidateConcat]
        {\inp xs \gennedby r_1, \out & \dinp ys \gennedby r_2, \dout}
        {\inp xs \append ys \gennedby r_1 \concat r_2, \dout}$
    $\inference
        [\ValidateEmpty]
        {}
        {\inp \nil \gennedby \nil, \out}$

    % literals
    $\inference
        [\ValidateLiteral]
        {}
        {\inp c \gennedby <c>, \env}$
    $\inference
        [\ValidateApplication]
        {\env_0 = \env
            & args' = \forall 1 \leq i \leq n. \; (\env_{i - 1}, \pathvar) arg'_i \gennedby arg_i, \env_i}
         {\inp <f \; args> \gens f(args'), \out}$

    % variables
    $\inference
        [\ValidateVarOne]
        {\envget{\lab} = x}
        {\inp x \gennedby \lab, \out}$
    $\inference
        [\ValidateVarTwo]
        {\envhasnt{\lab} & \arbitrary{x}}
        {\inp x \gennedby \lab, \outw{\lab}{x}}$

    $\inference
        [\ValidateUnionSolveLeft]
        {\envhasnt{\lab} & \inp x \gennedby r_1, \out}
        {\inp x \gennedby r_1 |_{\lab} r_2, \outw{\lab}{0}}$
    $\inference
        [\ValidateUnionSolveRight]
        {\envhasnt{\lab} & \inp x \gennedby r_2, \out}
        {\inp x \gennedby r_1 |_{\lab} r_2, \outw{\lab}{1}}$
    $\inference
        [\ValidateUnionLeft]
        {\envget{\lab} \mod 2 = 0 & \inp x \gennedby r_1, \out}
        {\inp x \gennedby r_1 |_{\lab} r_2, \out}$
    $\inference
        [\ValidateUnionRight]
        {\envget{\lab} \mod 2 = 1 & \inp x \gennedby r_2, \out}
        {\inp x \gennedby r_1 |_{\lab} r_2, \out}$
    
    $\inference
        [\ValidateStarSolve]
        {\envhasnt{\lab} & \exists n. \inpwe{\lab}{n} xs \gennedby r*_{\lab}, \out}
        {\inp xs \gennedby r*_{lab}, \out}$
    $\inference
         [\ValidateStar]
         { \envget{\lab} = n
            & \env_0 = \env
            & xs = \forall 1 \leq i \leq n . (\env_{i - 1}, \pathput{\lab}{i}) xs_i \gennedby r, \env_i
         }
         {\inp \foldconcat{xs} \gennedby r*_{\lab}, \env_n}$

    \caption{Validation Rules}
    \label{fig:validation-rules}
\end{center}
\end{figure*}

Once we have a run a function on randomly generated inputs,
we must \textit{validate} that the function's output
matches the output symbolic expression.
We define a relation
$\inp xs \gennedby e \gennedarrow ys, \out$,
such that $xs$ is a concrete list and $e$ is a symbolic expression
that we try to match to a \textit{prefix} of $xs$.
$ys$ is then the \textit{suffix} of $xs$ that was not matched to $e$.
Another way of looking at this
is that
$\inp xs \gennedby e \gennedarrow ys, \out$
says that there is some $xs'$ such that $xs = xs' :++ ys$,
and $xs'$ matches the symbolic expression $e$.
This allows, ``chaining'' together validation steps:
to check if some concrete list $xs$ matches a regular
expression $(r_1 \concat r_2) \concat r_3$,
we would first check if, for some $ys$
$\inp xs \gennedby r_1 \concat r_2 \gennedarrow ys, \out$,
and then pass $ys$ to $r_3$,
$(\env', \pathvar) ys \gennedby r_3 \gennedarrow zs, \env''$.
We would then know the match on $(r_1 \concat r_2) \concat r_3$ succeeded as a whole
if and only if $zs$ is the empty list--
\textit{every value} of $xs$ was matched to the symbolic expression,
with no suffix left over.


When validating,
we do not allow variables bindings to have constraints other than \text{True};
besides this, we construct the rules such that a list will be validated against a symbolic expression
if and only if the list could be generated from the symbolic expression.

Figure~\ref{fig:validation-rules} formalizes the validation rules.
Note that the actual operation of the validation rules is quite different from that of the generation rules.
During generation,
we randomly generate values at each binding and make use of those values later in constructing
a list.
During validation, we instead search for \textit{any} instantiation of bound variables
that will allow the match to succeed.
This requires delaying instantiation of variables until they are \textit{used},
and can be determined to correspond to some concrete expression.
In some cases, we must explore multiple possible matches
(for example, both branches of a union.)
This in reflected in the fact that the rules are \textit{not} entirely syntax-directed--
in the case that multiple rules apply,
the validator tries applying all of them,
in sequence,
until either (1) it finds a rule application that allows the match to succeed
or, (2) all possible rules, and thus the test as a whole, fail.


\mypara{List Constructors}
\ValidateCons, \ValidateConcat, and \ValidateEmpty
structurally match on a symbolic expression and a concrete list.
\ValidateCons
works very similarly to the corresponding generation rules,
\RuleCons, requiring the head and tail of the concrete list to match
against a head and tail in the symbolic expression.
Note that we require the head to \textit{completely} match
(the leftover suffix must be the empty list),
while the tail may pass on some elements to following pieces of the symbolic expression.
\ValidateEmpty allows an empty symbolic expression to match on any list,
returning the entire list as a suffix.
\bill{Why? We might want to actually show an inference tree}
\ValidateConcat ``non-deterministically'' splits the list $xs$
to match to each of the two sub-expressions in $r_1 \concat r_2$,
passing the suffix obtained from $r_1$ to the expression for $r_2$.
In practice, this non-determinism corresponds to an implementation sequentially (or in parallel)
trying every possible point to split.
While this might seem potentially expensive,
in practice most split points can be quickly discarded based on matching failures in the sub-expressions.

\mypara{Literals}
\ValidateLiteral matches a literal $c$ in the list to that same literal $c$ in the symbolic expression.
As with \ValidateCons and \ValidateEmpty,
this is not notably different then the corresponding generation rule,
\RuleLiteral

\mypara{Assignment and Variables}
\ValidateVarOne and \ValidateVarTwo
handle matching labels, $v \gennedby \lab$.
\ValidateVarOne applies if the label being matched
is contained is already in the environment--
in this case,
the rule checks that $\envget{\label}$ is indeed $v$.
\ValidateVarTwo allows a label that is not the environment
to match \textit{any} value,
and inserts that value in the environment to ensure that
mapping of the label is consistent
(i.e. if the variable appears elsewhere in the expression, then \ValidateVarOne will be applied.)

\ValidateAssignment
essentially treats assignment
$(\lab \Longleftarrow \mathit{True}) \mapsto e$
as a no-op,
simply directly continuing
matching on $e$.
This is because we want to put off assigning variables
in order to better align them to the concrete expression via, i.e. \ValidateVarTwo.
Note that assignments do still have an impact on
the meaning of an expression because they impact the computation of the dependency map--
thus, we do not omit assignments altogether. 

\mypara{Union}
We have four rules to handle a union $r_1 |_\lab r_2$.
Two of these rules,
\ValidateUnionSolveLeft and \ValidateUnionSolveRight,
handle the case where $\lab$ is not bound in the environment.
\ValidateUnionSolveLeft tries to match
the list against the left side of the union $r_1$,
binding $(\pathvar, \lab)$ to $0$ in the environment.
Correspondingly,
\ValidateUnionSolveRight tries to match against $r_2$,
binding $(\pathvar, \lab)$ to $1$.
\ValidateUnionLeft and \ValidateUnionRight
handle the case where $\lab$ is bound in
the environment.
In this case, no branching will occur,
as $\envget{\lab}$
will determine which of the two rules is applicable. 

\mypara{Kleene Closure}
Much like for the union,
we have seperate rules for kleene closures $r*_\lab$
depending on whether \lab is bound in the environment.
\ValidateStarSolve handles the case where
the label is not bound.
Given a concrete list $xs$
and the symoblic expression $r*_\lab$,
the rule
allows matching $r$ against $x$
some arbitrary number $n$ times,
i.e. setting $\env_0 = \env$,
the match is computed as
$\forall 1 \leq i < n . (\env_{i}, \pathput{\lab}{i}) ys_i \gennedby r \gennedarrow ys_{i + 1}, \env_{i + 1}$.
This roughly corresponds to a loop in which
(1) $xs$ is matched to $r$, obtaining some unmatched suffix $ys$,
and then (2) $xs$ is assigned the value $ys$, allowing another match to occur.
The system branches to explore each possible value of
$n$.
\ValidateStar handles matches on stars $*_\lab$ when \lab is mapped in the environment.
This works essentially the same way as in \ValidateStarSolve,
except that no branching is needed, as the value of $n$ is known.

\mypara{Apply}
\ValidateApplication
validates data constructor applications.
The rule works by checking that each
concrete argument matches each symbolic expression \textit{completely};
that is, we require that there be no left over suffix,
as there are no sub-expressions following the arguments to pass this suffix onto.
(Note that a use of the \ValidateApplication rule is typically
going to occur as either the root of a tree,
or a subtree following directly from matching a list head in \ValidateCons.)

% \bill{Somewhere note that predicates not usable in output}

% \par Validation happens slightly different than generation.
% Firstly, there are no constraints on variables during the validation phase.
% The purpose of validation is, given a symbolic expression and a list, to determine if
% it is possible (following the rules of generation) that the symbolic expression could
% generate the list (i.e $\forall e, xs. e \gens xs \Longleftrightarrow xs \gennedby e$).
% When we come across a label that has no value, we must try every value that label could
% take and check if one of them lets the output match the expression. If it can, we say the
% list was validated by the expression.



